Un módulo industrial consta de una placa metálica y una carcasa en contacto. Ambas pueden almacenar calor y cederlo al aire. Un calefactor actúa sobre la placa; al apagarlo, una prueba sintética registra \(80\,{}^\circ\mathrm C\) en la placa y \(20\,{}^\circ\mathrm C\) en la carcasa, en aire a \(20\,{}^\circ\mathrm C\). Después sólo se registra un termómetro adherido a la carcasa. No hay cambios de fase. Almacenar una energía \(C\,\Delta T\) eleva la temperatura de una parte en \(\Delta T\), si su capacidad térmica \(C\) se considera constante.

Proponga un modelo diferencial que permita interpretar ese registro. Decida qué temperaturas necesita representar, qué intercambios conserva y cómo los describe. Declare las condiciones iniciales y los parámetros que habría que conocer. Resuelva un caso sencillo que usted justifique, compruebe su comportamiento límite y explique cuándo sería aceptable reemplazarlo por un modelo más simple. Señale una observación que ayude a decidir entre ambas descripciones.
